\documentclass[11pt]{article}

\usepackage[utf8]{inputenc}
\usepackage[T1]{fontenc}
\usepackage{amsmath}
\usepackage{amssymb}
\usepackage{amsthm}
\usepackage[margin=1in]{geometry}
\usepackage{xcolor}
\usepackage{listings}
\usepackage{hyperref}

% ---- Lean listings: colour classes are (1) commands, (2) sorts and types,
% ---- (3) the names this development defines.
\lstdefinelanguage{Lean}{
  morekeywords=[1]{inductive,def,theorem,lemma,abbrev,structure,where,deriving,match,with,fun,if,then,else,do,let,in,instance,namespace,end,simp,induction,cases,rw,exact,intro,by,private,open,variable,section,set_option,return,mutual},
  morekeywords=[2]{Prop,Type,Sort,List,Nat,Bool,Option,Sum,Prod,Unit,UInt32,UInt64,Int32,Var,Ctor},
  morekeywords=[3]{Expr,Pure,Value,Root,Step,Ectx,Lower,Occs,Fits,TEnv,Multi,CHeap,Cell,Counted,Bound,Couple,Owns,TStep,TCtx,Ref,update,expand,uses,Closed,lookup,expand,Release},
  sensitive=true,
  morecomment=[l]{--},
  morecomment=[n]{/-}{-/},
  morestring=[b]",
  literate=
    {→}{{$\to$}}1 {←}{{$\leftarrow$}}1 {↔}{{$\leftrightarrow$}}1
    {∀}{{$\forall$}}1 {∃}{{$\exists$}}1 {∈}{{$\in$}}1 {∉}{{$\notin$}}1
    {≠}{{$\neq$}}1 {≤}{{$\leq$}}1 {≥}{{$\geq$}}1 {∧}{{$\wedge$}}1 {∨}{{$\vee$}}1
    {¬}{{$\neg$}}1 {⊆}{{$\subseteq$}}1 {×}{{$\times$}}1 {∘}{{$\circ$}}1
    {↦}{{$\mapsto$}}1 {≃}{{$\simeq$}}1 {⊥}{{$\bot$}}1 {∅}{{$\emptyset$}}1
    {⟨}{{$\langle$}}1 {⟩}{{$\rangle$}}1 {·}{{$\cdot$}}1
    {ρ}{{$\rho$}}1 {Γ}{{$\Gamma$}}1 {Δ}{{$\Delta$}}1 {α}{{$\alpha$}}1
    {β}{{$\beta$}}1 {σ}{{$\sigma$}}1 {τ}{{$\tau$}}1 {φ}{{$\varphi$}}1
    {₁}{{$_1$}}1 {₂}{{$_2$}}1 {₀}{{$_0$}}1 {′}{{$'$}}1
}
\lstset{
  language=Lean,
  inputencoding=utf8,
  extendedchars=true,
  basicstyle=\ttfamily\small,
  keywordstyle=[1]\color{blue!75!black}\bfseries,
  keywordstyle=[2]\color{teal!70!black},
  keywordstyle=[3]\color{violet!80!black},
  commentstyle=\color{black!45}\itshape,
  stringstyle=\color{orange!60!black},
  keepspaces=true,
  columns=fullflexible,
  showstringspaces=false,
  breaklines=true,
  frame=single,
  framerule=0.2pt,
  rulecolor=\color{black!25},
  xleftmargin=1em,
  aboveskip=0.8em, belowskip=0.8em,
}
\newcommand{\code}[1]{\texttt{\detokenize{#1}}}
\newcommand{\rulename}[1]{\textsc{#1}}
\newcommand{\lref}[1]{\code{#1}}

\newtheorem{theorem}{Theorem}
\newtheorem{lemma}[theorem]{Lemma}
\theoremstyle{remark}
\newtheorem{remark}[theorem]{Remark}

\title{Reference count one at the write}
\author{Gavin Mendel-Gleason}
\date{}

\begin{document}
\maketitle

\begin{abstract}
A reference-counted runtime updates an object in place exactly when its count is
one. Implementations decide this at run time, because one compiled body serves
owning and borrowing call sites; but the count is determined by which references
the term and the surrounding structure hold, so it is a static property of the
program. This paper formalises a pure core, a reference-counted target over a
counted heap, and the lowering between them, and proves the pair that
characterises the target's write: an owned store is destructive and keeps its
location, a shared store copies. Everything is mechanised in Lean~4 with no
library dependency. The lowering is proved total and deterministic, not correct;
that is stated.
\end{abstract}

\section{The guarantee}

\noindent The write is destructive when its operand is owned once and nothing else
names the cell; otherwise it copies. That is the pair proved here, and it is the
only claim made about the target's write.

\begin{lstlisting}
theorem store_owned_is_destructive {c : CHeap} {hs : List Ref} {ρ : TEnv} {m : Multi}
    {y : Var} {r : Ref} {i : Nat} {v : Expr} {ks : List Ref} (hc : Couple c hs ρ m)
    (hρ : ρ.lookup y = some r) (hown : Owns ρ m (.store (.var y) i v))
    (hm : (expand m).count r = 1) (hslot : (CHeap.slots c).count r = 0)
    (hrun : hs.count r = 0) (hks : r ∉ ks) :
    (update c (hs ++ expand m) r ks).1.lookup r = some ks

theorem store_shared_copies {c : CHeap} {hs : List Ref} {ρ : TEnv} {m : Multi}
    {y : Var} {r : Ref} {i : Nat} {v : Expr} {ks : List Ref} (hc : Couple c hs ρ m)
    (hρ : ρ.lookup y = some r) (hown : Owns ρ m (.store (.var y) i v))
    (hm : 2 ≤ (expand m).count r ∨ 1 ≤ (CHeap.slots c).count r) (hks : r ∉ ks) :
    (update c (hs ++ expand m) r ks).1.lookup (fresh c ks) = some ks ∧
      (update c (hs ++ expand m) r ks).1.lookup r = CHeap.lookup c r ∧
      (update c (hs ++ expand m) r ks).1.count r = CHeap.count c r - 1
\end{lstlisting}

\noindent \lref{Couple} is the coupling invariant (\S\ref{sec:account}), \lref{Owns}
is the term-indexed account of the term's own references, \lref{expand m} is the
reference list the account stands for, and \lref{update} is the heap's two-branch
write (\S\ref{sec:heap}). The two hypotheses the account cannot supply --- no slot
and no run handle naming the cell --- are named rather than hidden
(\S\ref{sec:pair}). Read together the two theorems say: the write is in place
exactly where the account counts the operand once.

\section{The counted heap}
\label{sec:heap}

\noindent A cell holds the references its slots name and the count the runtime keeps
for it; a heap lists the live cells.

\begin{lstlisting}
structure Cell where
  slots : List Ref
  count : Nat
deriving DecidableEq, Repr

abbrev CHeap := List (Ref × Cell)
\end{lstlisting}

\noindent One reference names one cell, so the count a heap keeps at a reference is
not a choice among cells.

\begin{lstlisting}
def CHeap.WF (c : CHeap) : Prop := (CHeap.refs c).Nodup
\end{lstlisting}

\noindent No live cell has count zero, because a release deallocates at zero.

\begin{lstlisting}
def CHeap.Live (c : CHeap) : Prop := c.find? (fun e => e.2.count == 0) = none
\end{lstlisting}

\noindent A cell's count is what it must be: its occurrences among the live cells'
slots, plus the handles held outside the heap --- the references the running term
holds.

\begin{lstlisting}
def Counted (c : CHeap) (handles : List Ref) : Prop :=
  ∀ r, CHeap.count c r = refCount (CHeap.presentation c) handles r
\end{lstlisting}

\noindent The write branches on that count: at one the cell keeps its location (and
the references its old slots held are dropped), otherwise a fresh cell takes the
written slots and the old cell loses a reference.

\begin{lstlisting}
def update (c : CHeap) (handles : List Ref) (r : Ref) (ks : List Ref) : CHeap × List Ref :=
  if CHeap.count c r = 1 then
    (CHeap.cascade (CHeap.writeSlots c r ks), handles)
  else
    ((fresh c ks, cellOf ks) :: CHeap.addOccurrences (CHeap.adjust c r (· - 1)) ks,
      fresh c ks :: dropHandle handles r)
\end{lstlisting}

\noindent The two branches, as already proved of the heap: a destructive write is
read back at the same location, a copying write at a location the heap did not list.

\begin{lstlisting}
theorem update_preserves_location_of_exclusive {c : CHeap} {handles : List Ref}
    (hW : CHeap.WF c) (hc : Counted c handles) {r : Ref} {ks : List Ref}
    (h₁ : Exclusive (CHeap.presentation c) handles r) (hhandle : 1 ≤ handles.count r)
    (hks : r ∉ ks) :
    (update c handles r ks).1.lookup r = some ks

theorem update_copies_of_shared {c : CHeap} {handles : List Ref} (hc : Counted c handles)
    {r : Ref} {ks : List Ref} (h₂ : ¬ Exclusive (CHeap.presentation c) handles r)
    (hhandle : 1 ≤ handles.count r) (hks : r ∉ ks) :
    (update c handles r ks).1.lookup (fresh c ks) = some ks ∧
      (update c handles r ks).1.lookup r = CHeap.lookup c r ∧
      (update c handles r ks).1.count r = CHeap.count c r - 1 ∧
      fresh c ks ∉ CHeap.refs c
\end{lstlisting}

\section{The source calculus}
\label{sec:source}

\noindent One term type serves both calculi. Nodes are nullary, unary or binary; an
$n$-ary node is a spine. The write is a subject, a slot index and a value.

\begin{lstlisting}
inductive Expr where
  | var (x : Var)
  | lam (x : Var) (body : Expr)
  | app (f a : Expr)
  | bind (x : Var) (e₁ e₂ : Expr)
  | matchC (s : Expr) (c : Ctor) (ys : List Var) (body : Expr)
  | dup (x : Var) (e : Expr)
  | drop (x : Var) (e : Expr)
  | con0 (c : Ctor)
  | con1 (c : Ctor) (a : Expr)
  | con2 (c : Ctor) (a b : Expr)
  | store (s : Expr) (i : Nat) (v : Expr)
deriving DecidableEq, Repr
\end{lstlisting}

\noindent The source is the fragment with neither instruction; the free-variable
function and the closedness discipline are what make the substitution below sound.

\begin{lstlisting}
inductive Pure : Expr → Prop where
  | var (x : Var) : Pure (.var x)
  | lam (x : Var) (body : Expr) : Pure body → Pure (.lam x body)
  | app (f a : Expr) : Pure f → Pure a → Pure (.app f a)
  | bind (x : Var) (e₁ e₂ : Expr) : Pure e₁ → Pure e₂ → Pure (.bind x e₁ e₂)
  | matchC (s : Expr) (c : Ctor) (ys : List Var) (body : Expr) :
      Pure s → Pure body → Pure (.matchC s c ys body)
  | con0 (c : Ctor) : Pure (.con0 c)
  | con1 (c : Ctor) (a : Expr) : Pure a → Pure (.con1 c a)
  | con2 (c : Ctor) (a b : Expr) : Pure a → Pure b → Pure (.con2 c a b)
  | store (s : Expr) (i : Nat) (v : Expr) : Pure s → Pure v → Pure (.store s i v)

def uses : Expr → List Var
  | .var x => [x]
  | .lam x body => (uses body).filter (fun y => y ≠ x)
  | .app f a => uses f ++ uses a
  | .bind x e₁ e₂ => uses e₁ ++ (uses e₂).filter (fun y => y ≠ x)
  | .matchC s _ ys body => uses s ++ (uses body).filter (fun y => y ∉ ys)
  | .dup x e => x :: uses e
  | .drop x e => x :: uses e
  | .con0 _ => []
  | .con1 _ a => uses a
  | .con2 _ a b => uses a ++ uses b
  | .store s _ v => uses s ++ uses v

def Closed (e : Expr) : Prop := uses e = []
\end{lstlisting}

\noindent Values are the canonical forms the reduction's rules match on.

\begin{lstlisting}
inductive Value : Expr → Prop where
  | var (x : Var) : Value (.var x)
  | lam (x : Var) (body : Expr) : Value (.lam x body)
  | con0 (c : Ctor) : Value (.con0 c)
  | con1 {c : Ctor} {a : Expr} (h : Value a) : Value (.con1 c a)
  | con2 {c : Ctor} {a b : Expr} (h₁ : Value a) (h₂ : Value b) : Value (.con2 c a b)
\end{lstlisting}

\noindent The redexes: application of an abstraction, the two binders, a match on a
value node, and the write read as a functional update.

\begin{lstlisting}
inductive Root : Expr → Expr → Prop where
  | beta {x : Var} {body v : Expr} :
      Value v → Root (.app (.lam x body) v) (subst x v body)
  | bindVal {x : Var} {v body : Expr} :
      Value v → Root (.bind x v body) (subst x v body)
  | matchVal {s : Expr} {c : Ctor} {ys : List Var} {body : Expr} : Value s →
      nodeTag? s = some c → ys.length = (nodeFields s).length →
      Root (.matchC s c ys body) (substAll ys (nodeFields s) body)
  | storeVal {s : Expr} {i : Nat} {v n' : Expr} : Value s → Value v →
      slotUpdate s i v = some n' → Root (.store s i v) n'
\end{lstlisting}

\noindent The reduction is a redex, or a step inside one of the contexts below.
There is no rule inside a match's subject, a write's subject or a node's fields:
the source is written so that those are values already.

\begin{lstlisting}
inductive Step : Expr → Expr → Prop where
  | root {e e' : Expr} : Root e e' → Step e e'
  | appFun {f f' a : Expr} : Step f f' → Step (.app f a) (.app f' a)
  | appArg {f a a' : Expr} : Value f → Step a a' → Step (.app f a) (.app f a')
  | bindFun {e₁ e₁' : Expr} {x : Var} {e₂ : Expr} : Step e₁ e₁' →
      Step (.bind x e₁ e₂) (.bind x e₁' e₂)
\end{lstlisting}

\noindent The contexts the source steps in.

\begin{lstlisting}
inductive Ectx where
  | hole
  | appL (E : Ectx) (a : Expr)
  | appR (f : Expr) (E : Ectx)
  | bindL (x : Var) (E : Ectx) (e₂ : Expr)
\end{lstlisting}

\noindent Substitution is naive, and therefore capture free only on closed
arguments: this is the discipline the source is used at.

\begin{lstlisting}
theorem uses_subst_of_closed : ∀ (e : Expr) (x : Var) (v : Expr), Pure e → Closed v →
    uses (subst x v e) = (uses e).filter (fun y => y ≠ x)

theorem step_preserves_closed {e e' : Expr} (hs : Step e e') : Pure e → Closed e → Closed e'
\end{lstlisting}

\noindent The reduction is a function, which is what lets the lowering be one too.

\begin{lstlisting}
theorem step_deterministic : ∀ {e e₁ : Expr}, Step e e₁ → ∀ {e₂ : Expr}, Step e e₂ → e₁ = e₂
\end{lstlisting}

\section{The lowering}
\label{sec:lower}

\noindent A derivation relates a source term to its target under a borrowed and an
owned multiset of occurrences. Availability is about names, not multiplicities: a
name used twice may be fitting at one owned occurrence, and the strategy supplies
the second use with a duplication.

\begin{lstlisting}
def Fits (Δ Γ : Occs) (e : Expr) : Prop := ∀ x ∈ uses e, x ∈ Γ ∨ x ∈ Δ
\end{lstlisting}

\noindent The relation: one rule per source construct. The splits are the functions
\lref{claimEnv}, \lref{bindEnv} and \lref{armEnv} rather than an existential, and
each rule's conclusion wraps its translation in \lref{wrap}, the prologue that
duplicates and releases.

\begin{lstlisting}
inductive Lower : Occs → Occs → Expr → Expr → Prop where
  | var {Δ Γ x} (h : x ∈ Γ ∨ x ∈ Δ) :
      Lower Δ Γ (.var x) (wrap Γ [x] (.var x))
  | lam {Δ Γ x body body'} (h : Lower (forget Δ x) (bindEnv Γ x (uses body)) body body') :
      Lower Δ Γ (.lam x body) (wrap Γ (uses (.lam x body)) (.lam x body'))
  | app {Δ Γ f a f' a'} (hf : Lower Δ (claimEnv Γ (uses f)) f f')
      (ha : Lower Δ (claimEnv Γ (uses a)) a a') :
      Lower Δ Γ (.app f a) (wrap Γ (uses (.app f a)) (.app f' a'))
  | bind {Δ Γ x e₁ e₂ e₁' e₂'} (h₁ : Lower Δ (claimEnv Γ (uses e₁)) e₁ e₁')
      (h₂ : Lower (forget Δ x) (bindEnv Γ x (uses e₂)) e₂ e₂') :
      Lower Δ Γ (.bind x e₁ e₂) (wrap Γ (uses (.bind x e₁ e₂)) (.bind x e₁' e₂'))
  | matchC {Δ Γ s c ys body s' body'} (hs : Lower Δ (claimEnv Γ (uses s)) s s')
      (h : Lower (forgetAll Δ ys) (armEnv Γ ys (uses body)) body body') :
      Lower Δ Γ (.matchC s c ys body)
        (wrap Γ (uses (.matchC s c ys body)) (.matchC s' c ys body'))
  | con0 {Δ Γ c} : Lower Δ Γ (.con0 c) (wrap Γ [] (.con0 c))
  | con1 {Δ Γ c a a'} (h : Lower Δ (claimEnv Γ (uses a)) a a') :
      Lower Δ Γ (.con1 c a) (wrap Γ (uses (.con1 c a)) (.con1 c a'))
  | con2 {Δ Γ c a b a' b'} (ha : Lower Δ (claimEnv Γ (uses a)) a a')
      (hb : Lower Δ (claimEnv Γ (uses b)) b b') :
      Lower Δ Γ (.con2 c a b) (wrap Γ (uses (.con2 c a b)) (.con2 c a' b'))
  | store {Δ Γ s i v s' v'} (hs : Lower Δ (claimEnv Γ (uses s)) s s')
      (hv : Lower Δ (claimEnv Γ (uses v)) v v') :
      Lower Δ Γ (.store s i v)
        (wrap Γ (uses (.store s i v)) (.store s' i v'))
\end{lstlisting}

\noindent The relation is exactly the graph of the computable strategy \lref{lowerB}
--- \lref{lowerB_sound} and \lref{lowerB_complete} are its two directions --- so
the strategy's own properties are the relation's. It is total on the fragment and
deterministic, and the second is why the splits are canonical: with an arbitrary
split of a binding's environment the one required duplication could be placed in
either subterm and the targets would differ.

\begin{lstlisting}
theorem lower_total {Δ Γ : Occs} {e : Expr} (hp : Pure e) (hfit : Fits Δ Γ e) :
    ∃ e', Lower Δ Γ e e'

theorem lower_deterministic {Δ Γ : Occs} {e e₁ e₂ : Expr} (h₁ : Lower Δ Γ e e₁)
    (h₂ : Lower Δ Γ e e₂) : e₁ = e₂
\end{lstlisting}

\noindent These are properties of the translation. Nothing here says the target
computes what the source computes; that is \S\ref{sec:future}.

\section{The target and its runtime}
\label{sec:target}

\noindent A configuration carries the heap, the handles the run holds, an
environment from variables to references, the term's own account of its references,
and the term.

\begin{lstlisting}
abbrev TEnv := List (Var × Ref)
abbrev Multi := List (Var × Ref)
def expand (m : Multi) : List Ref := m.map (fun p => p.2)
\end{lstlisting}

\noindent The account is one pair per occurrence the term itself holds. Each
instruction moves it in step with the heap: a duplication increments the count and
adds an occurrence; a release decrements, or deallocates at a count of one; a
binding or application of a variable moves an occurrence between names; a match
binds the arm's names to the cell's slots; and the write is the heap's own
\lref{update}, given the account as outside handles in addition to the run's.

\begin{lstlisting}
inductive TStep :
    CHeap → List Ref → TEnv → Multi → Expr → CHeap → List Ref → TEnv → Multi → Expr → Prop where
  | dup {c : CHeap} {hs : List Ref} {ρ : TEnv} {m : Multi} {y : Var} {r : Ref} {e : Expr}
      (hρ : ρ.lookup y = some r) (hlive : r ∈ CHeap.refs c) :
      TStep c hs ρ m (.dup y e) (dup c hs r).1 hs ρ (Multi.insert m y r) e
  | dropShared {c : CHeap} {hs : List Ref} {ρ : TEnv} {m : Multi} {y : Var} {r : Ref}
      {e : Expr} (hρ : ρ.lookup y = some r) (h₂ : 2 ≤ CHeap.count c r)
      (hsub : (y, r) ∈ m) :
      TStep c hs ρ m (.drop y e) (decRef c hs r).1 hs ρ (Multi.dropOcc m y r) e
  | dropFinal {c : CHeap} {hs : List Ref} {ρ : TEnv} {m : Multi} {y : Var} {r : Ref}
      {e : Expr} (hρ : ρ.lookup y = some r) (h₁ : CHeap.count c r = 1)
      (hsub : (y, r) ∈ m) :
      TStep c hs ρ m (.drop y e) (release c hs r).1 hs ρ (Multi.dropOcc m y r) e
  | bindVar {c : CHeap} {hs : List Ref} {ρ : TEnv} {m : Multi} {x z : Var} {e₂ : Expr}
      {r : Ref} (hρ : ρ.lookup z = some r) (hsub : (z, r) ∈ m)
      (hhigh : nameOnly m x r = true) :
      TStep c hs ρ m (.bind x (.var z) e₂) c hs (TEnv.set ρ x r)
        (Multi.insert (Multi.dropOcc m z r) x r) e₂
  | bindVal {c : CHeap} {hs : List Ref} {ρ : TEnv} {m : Multi} {x : Var} {v e₂ : Expr}
      {c' : CHeap} {hs' : List Ref} {r : Ref} (hv : ValRef c hs ρ v c' hs' r)
      (halloc : ∀ z, v ≠ .var z) (hfree : freeName m x = true) :
      TStep c hs ρ m (.bind x v e₂) c' (dropHandle hs' r) (TEnv.set ρ x r)
        (Multi.insert m x r) e₂
  | appVar {c : CHeap} {hs : List Ref} {ρ : TEnv} {m : Multi} {x z : Var} {body : Expr}
      {r : Ref} (hρ : ρ.lookup z = some r) (hsub : (z, r) ∈ m)
      (hhigh : nameOnly m x r = true) :
      TStep c hs ρ m (.app (.lam x body) (.var z)) c hs (TEnv.set ρ x r)
        (Multi.insert (Multi.dropOcc m z r) x r) body
  | appVal {c : CHeap} {hs : List Ref} {ρ : TEnv} {m : Multi} {x : Var} {body v : Expr}
      {c' : CHeap} {hs' : List Ref} {r : Ref} (hv : ValRef c hs ρ v c' hs' r)
      (halloc : ∀ z, v ≠ .var z) (hfree : freeName m x = true) :
      TStep c hs ρ m (.app (.lam x body) v) c' (dropHandle hs' r) (TEnv.set ρ x r)
        (Multi.insert m x r) body
  | matchC {c : CHeap} {hs : List Ref} {ρ : TEnv} {m : Multi} {z : Var} {ctag : Ctor}
      {ys : List Var} {body : Expr} {r : Ref} (hρ : ρ.lookup z = some r)
      (hlen : ys.length = (CHeap.slotsAt c r).length) (hfree : ∀ p ∈ m, p.1 ∉ ys) :
      TStep c hs ρ m (.matchC (.var z) ctag ys body) c hs
        (TEnv.bindAll ρ ys (CHeap.slotsAt c r)) m body
  | storeExclusive {c : CHeap} {hs : List Ref} {ρ : TEnv} {m : Multi} {y : Var} {i : Nat}
      {v : Expr} {c1 : CHeap} {hs1 : List Ref} {wv r : Ref} (hρ : ρ.lookup y = some r)
      (hv : ValRef c hs ρ v c1 hs1 wv) (hi : i < (CHeap.slotsAt c1 r).length)
      (hchk : writeCheck c1 r (storeSlots c1 r i wv) = true) (h₁ : CHeap.count c1 r = 1)
      (hsub : (y, r) ∈ m) :
      TStep c hs ρ m (.store (.var y) i v)
        (update c1 (hs1 ++ expand m) r (storeSlots c1 r i wv)).1 hs1 ρ m (.var y)
  | storeShared {c : CHeap} {hs : List Ref} {ρ : TEnv} {m : Multi} {y : Var} {i : Nat}
      {v : Expr} {c1 : CHeap} {hs1 : List Ref} {wv r : Ref} (hρ : ρ.lookup y = some r)
      (hv : ValRef c hs ρ v c1 hs1 wv) (hi : i < (CHeap.slotsAt c1 r).length)
      (hchk : writeCheck c1 r (storeSlots c1 r i wv) = true) (hsub : (y, r) ∈ m)
      (hgone : nameGone m y r = true) (h₂ : CHeap.count c1 r ≠ 1) :
      TStep c hs ρ m (.store (.var y) i v)
        (update c1 (hs1 ++ expand m) r (storeSlots c1 r i wv)).1 hs1
        (TEnv.set ρ y (fresh c1 (storeSlots c1 r i wv)))
\end{lstlisting}

\noindent The two store rules are the two branches of \lref{update}, and the only
difference between them is the count. Each reads the subject's occurrence out of the
term's own account (\lref{hsub}), so an account that does not hold it cannot step.

\noindent The positions the runtime steps in: a step inside a well-formed context is
a step. The rules above are the steps at a configuration's root; the remaining rules
of \lref{TStep} --- the function position of an application, its argument while the
function is a value, a binding's right-hand side, and a node's fields --- are these
contexts one by one, which \lref{tstep_of_plug} states.

\begin{lstlisting}
inductive TCtx where
  | hole
  | appL (E : TCtx) (a : Expr)
  | appR (f : Expr) (E : TCtx)
  | bindL (x : Var) (E : TCtx) (e₂ : Expr)
  | con1L (ct : Ctor) (E : TCtx)
  | con2L (ct : Ctor) (E : TCtx) (b : Expr)
  | con2R (ct : Ctor) (a : Expr) (E : TCtx)

theorem tstep_of_plug : ∀ {E : TCtx} {c : CHeap} {hs : List Ref} {ρ : TEnv} {m : Multi}
    {r e' : Expr} {c' : CHeap} {hs' : List Ref} {ρ' : TEnv} {m' : Multi},
    TCtx.WF E → TStep c hs ρ m r c' hs' ρ' m' e' →
      TStep c hs ρ m (E.plug r) c' hs' ρ' m' (E.plug e')
\end{lstlisting}

\section{The account, and the pair}
\label{sec:account}

\noindent Exact tagging: each pair names the reference its variable is bound to, not
merely some reference. The weaker reading admits a mis-tagged pair and makes the
theorems below arithmetic.

\begin{lstlisting}
def Bound (ρ : TEnv) (m : Multi) : Prop := ∀ p ∈ m, ρ.lookup p.1 = some p.2
\end{lstlisting}

\noindent The account is of the term: each consuming instruction reads its
occurrence out of it, and a store reads the occurrence its subject names.

\begin{lstlisting}
inductive Owns (ρ : TEnv) : Multi → Expr → Prop where
  | drop {m : Multi} {y : Var} {r : Ref} {e : Expr} (hρ : ρ.lookup y = some r)
      (hmem : (y, r) ∈ m) : Owns ρ m (.drop y e)
  | bindVar {m : Multi} {x z : Var} {r : Ref} {e₂ : Expr} (hρ : ρ.lookup z = some r)
      (hmem : (z, r) ∈ m) (hheld : ∀ p ∈ m, p.1 = x → p.2 = r) :
      Owns ρ m (.bind x (.var z) e₂)
  | appVar {m : Multi} {x z : Var} {r : Ref} {body : Expr} (hρ : ρ.lookup z = some r)
      (hmem : (z, r) ∈ m) (hheld : ∀ p ∈ m, p.1 = x → p.2 = r) :
      Owns ρ m (.app (.lam x body) (.var z))
  | store {m : Multi} {y : Var} {r : Ref} {i : Nat} {v : Expr} (hρ : ρ.lookup y = some r)
      (hmem : (y, r) ∈ m) : Owns ρ m (.store (.var y) i v)
\end{lstlisting}

\noindent The coupling invariant: the term's own references are handles. That is the
whole content of the middle conjunct.

\begin{lstlisting}
def Couple (c : CHeap) (hs : List Ref) (ρ : TEnv) (m : Multi) : Prop :=
  CHeap.WF c ∧ Counted c (hs ++ expand m) ∧ Bound ρ m
\end{lstlisting}

\noindent Every step preserves it, contexts included.

\begin{lstlisting}
theorem couple_step {c : CHeap} {hs : List Ref} {ρ : TEnv} {m : Multi} {e : Expr}
    {c' : CHeap} {hs' : List Ref} {ρ' : TEnv} {m' : Multi} {e' : Expr}
    (h : TStep c hs ρ m e c' hs' ρ' m' e') (hc : Couple c hs ρ m) :
    Couple c' hs' ρ' m'
\end{lstlisting}

\subsection{The two hypotheses, and where they come from}
\label{sec:pair}

\noindent \lref{store_owned_is_destructive} cannot derive, from the account, that no
heap slot and no run handle names the cell. Both are what a nested update's take-out
provides: the element is removed from its parent's slots, and the pair the take-out
hands back is destructured and its handle released before the focus is checked. The
mechanisation names them in the theorem and cites the take-out lemmas that supply
them.

\noindent The second theorem is what keeps the first from being vacuous: a subject a
second reference names \emph{must} copy, or the update would be visible through the
alias. A statement that every store is destructive is simply false, and the pair is
the honest form of the claim.

\section{Discussion}

\noindent \textbf{The guarantee is static although the runtime tests.} The count at
a write is the term's own occurrences plus the structure's: a function of the term,
the heap and the environment, and not a value read from a run. A test emitted by a
compiler is a residual of compiling one body for two ownership situations. A program
whose ownership is known does not need the test to be proved about it; only
executed.

\noindent \textbf{What is trusted.} Three things. The take-out's contribution: that
a nested update's cell has no slot and no handle naming it is a property of the
nesting, and is taken here as the two named hypotheses. The compiler's lowering:
that a compiler turns a language's write paths into the shape modelled here is not
proved; in the accompanying repository that link is \emph{checked} by reading the
post-reference-counting intermediate representation the pinned toolchain emits and
asserting the shape on one executing path, failing the build when it does not hold.
The runtime's primitive: the branch that stores in place is a runtime test whose
meaning is documented in the runtime and not modelled here.

\noindent \textbf{What is not proved.} The lowering's semantic correctness: no
theorem says the target computes what the source computes, or that it is garbage
free. What is proved of the lowering is totality and determinism. The declarative
relation of the paper this follows is not the one formalised: the strategy's splits
are canonical, for determinism, so the results are properties of \emph{this}
strategy, which is what a compiler implements.

\section{Mechanisation, and related work}
\label{sec:future}

\noindent The development is roughly 5000 lines of Lean~4 across five modules
against the core library only, with no mathematics library dependency. There is no
\code{sorry}, no evaluation by native decision, and every representative theorem's
axiom list is asserted, so a proof that acquired a foreign axiom fails the build.
Two design facts the formalisation forced are recorded as toolchain facts: terms are
a flat inductive, because a nested one cannot be inducted on or derive decidable
equality; and an evaluated check needs a structural recursion, because the kernel
does not reduce a well-founded one.

\noindent \textbf{Reference counting with reuse.} In-place update of a shared value,
and the precise counting with ownership that makes it possible, is the design of
Lean~4's runtime \cite{ullrich19counting} and is formalised as the linear resource
calculus $\lambda^1$ in the Pe\TeX{} paper \cite{reinking21perceus}, which proves
the counting sound and garbage free. This development follows that shape and
differs where this paper's subject differs: the deliverable is the count at the
write, as a pair, with the hypotheses the account cannot supply named; and the
translation is a canonical strategy rather than the declarative relation.

\noindent \textbf{Linearity, and formal models of counting.} An ordinary type system
does not track uniqueness: Linear Haskell \cite{bernardy18linear} gives a linear
arrow, and its guarantee is about the use of a value, not about a count in a heap.
A formal model of reference counting for code generation appears in F\'erey and
Shankar \cite{ferrey16codegen}. Neither is a count-one guarantee for a write.

\noindent \textbf{Future work.} Three steps: the lowering's semantic correctness
(result agreement and garbage freedom); the surface encoding that would make ``every
write in an accepted program is destructive'' a statement about a client's programs
rather than about a strategy; and a verified account, or a soundness theorem for the
check, of the compiler's pass.

\begin{thebibliography}{9}

\bibitem{ullrich19counting}
Sebastian Ullrich and Leonardo de Moura.
\newblock Counting immutable beans: reference counting optimized for purely
functional programming.
\newblock In \emph{Proceedings of the 31st Symposium on Implementation and
Application of Functional Languages (IFL '19)}, 2019. arXiv:1908.05647.

\bibitem{reinking21perceus}
Alex Reinking, Ningning Xie, Leonardo de Moura, and Daan Leijen.
\newblock Pe\TeX{}: garbage free reference counting with reuse.
\newblock In \emph{Proceedings of the 42nd ACM SIGPLAN International Conference on
Programming Language Design and Implementation (PLDI '21)}, 2021.
\newblock Extended version: Microsoft Research Technical Report MSR-TR-2020-42,
which contains the proofs.

\bibitem{bernardy18linear}
Jean-Philippe Bernardy, Mathieu Boespflug, Ryan R. Newton, Simon Peyton Jones, and
Arnaud Spiwack.
\newblock Linear Haskell: practical linearity in a higher-order polymorphic
language.
\newblock In \emph{Proceedings of the 45th ACM SIGPLAN Symposium on Principles of
Programming Languages (POPL '18)}, 2018.

\bibitem{ferrey16codegen}
Gaspard F\'erey and Natarajan Shankar.
\newblock Code generation using a formal model of reference counting.
\newblock In \emph{NASA Formal Methods (NFM 2016)}, pages 150--165, 2016.

\bibitem{demoura21lean4}
Leonardo de Moura and Sebastian Ullrich.
\newblock The Lean 4 theorem prover and programming language.
\newblock In \emph{Automated Deduction (CADE-28)}, 2021.

\bibitem{scidonia-lean}
Gavin Mendel-Gleason.
\newblock The mechanisation of this paper, in the Scidonia \code{lean} repository:
\code{Scidonia/RC/} and \code{Scidonia/FracPerm/StateVar/RCHeap.lean}, with the
checks in \code{ScidoniaTest/RC.lean} and \code{ScidoniaTest/Target.lean}.
\newblock \url{https://github.com/scidonia/lean-examples}.

\end{thebibliography}

\end{document}
